\section{Varying actual original-line data define intrinsic numerical FS functions}
\label{sec:native-numerical-fubini-study}

\begin{definition}[Whole original rational-line FS data]
\label{def:native-numerical-whole-section-fs-data}
\lean{EnvelopeProperty.NativeNumericalFubiniStudy.FSData,EnvelopeProperty.NativeNumericalFubiniStudy.FSData.numericalClass,EnvelopeProperty.NativeNumericalFubiniStudy.FSData.eval,EnvelopeProperty.NativeNumericalFubiniStudy.FSData.Nonzero,EnvelopeProperty.NativeNumericalFubiniStudy.FSData.IsRational}
\leanok
\uses{thm:native-real-numerical-classes,thm:scheme-global-section-weighted-norm-isometry}
\source{bj22-trivval:defi:FS,fhx-envelope-20261028:eq:FS-function}
On an actual original scheme, a datum consists of an actual invertible line
$M$, a positive integer $m$, and a finite family of pairs $(s,\lambda)$ of
actual whole global sections of $M$ and real weights. Its numerical class
on a projective original model is $m^{-1}[M]$. Its function at the same
original native points of a proper model is
\[
 m^{-1}\max_{(s,\lambda)}\bigl(\log\|s\|+\lambda\bigr),
\]
using the true original proper-model norm. Zero sections and common zeroes
are permitted, with $\log0=-\infty$. A datum is nonzero when at least one of
its actual sections is nonzero, and rational when every actual weight is
rational. The line is part of the actual data and can vary freely.
\end{definition}

\begin{theorem}[Actual nonzero sections exactly detect generic finiteness]
\label{thm:native-numerical-fs-nonvanishing}
\lean{EnvelopeProperty.NativeNumericalFubiniStudy.FSData.eval_ne_bot_iff,EnvelopeProperty.NativeNumericalFubiniStudy.FSData.nonzero_iff_exists_eval_ne_bot}
\leanok
\uses{def:native-numerical-whole-section-fs-data,thm:scheme-original-global-section-norm-detection}
At every actual original native point, the datum evaluates to a value
other than $-\infty$ precisely when one of its actual sections has nonzero
true norm there. If the proper original scheme is integral, the family is
nonzero if and only if its evaluation is not identically $-\infty$.
\end{theorem}
\begin{proof}
A weighted term is $-\infty$ exactly when its section norm vanishes, because
the weight is finite. A finite supremum equals $-\infty$ exactly when all
its terms do. The reciprocal denominator is strictly positive, so it
preserves this condition. This proves the pointwise assertion, including
empty families and common zeroes. On an integral proper original scheme,
an actual nonzero whole section has nonzero true norm at an actual original
native point by the genuine section-norm detection theorem. Such a term
makes the evaluation nontrivial. Conversely, a nontrivial evaluation yields
a term of nonzero norm, and the zero whole section has zero true norm.
\end{proof}

\begin{definition}[FS of an actual numerical class]
\label{def:native-numerical-whole-section-fs}
\lean{EnvelopeProperty.NativeNumericalFubiniStudy.IsFS,EnvelopeProperty.NativeNumericalFubiniStudy.IsRationalFS}
\leanok
\uses{def:native-numerical-whole-section-fs-data}
A function is numerical FS of class $\theta$ if it is the entire evaluation
of some nonzero actual datum whose actual class $m^{-1}[M]$ equals $\theta$.
The rational variant additionally requires rational weights. This permits
every actual original line and denominator; it has no fixed-line-ray
restriction and no supplied realization or comparison certificate.
\end{definition}

\begin{theorem}[Numerical FS retains its genuine original nontriviality]
\label{thm:native-numerical-fs-generically-finite}
\lean{EnvelopeProperty.NativeNumericalFubiniStudy.IsFS.exists_ne_bot,EnvelopeProperty.NativeNumericalFubiniStudy.IsRationalFS.isFS}
\leanok
\uses{def:native-numerical-whole-section-fs,thm:native-numerical-fs-nonvanishing}
On an integral projective original scheme, every numerical FS function is
not identically $-\infty$. Every rational numerical FS function is numerical
FS of the same class and the same entire original function.
\end{theorem}
\begin{proof}
Use the actual defining nonzero family and its preceding nonvanishing
criterion. Rational weights are a further property of the same datum;
forgetting this property retains its line, denominator, true class,
nonzero sections and entire evaluation.
\end{proof}

\begin{theorem}[Every existing original fixed-line datum has its genuine numerical class]
\label{thm:native-numerical-fs-from-ray}
\lean{EnvelopeProperty.NativeNumericalFubiniStudy.FSData.ofRay,EnvelopeProperty.NativeNumericalFubiniStudy.FSData.ofRay_eval,EnvelopeProperty.NativeNumericalFubiniStudy.FSData.ofRay_numericalClass,EnvelopeProperty.NativeNumericalFubiniStudy.isFS_of_ray}
\leanok
\uses{def:native-numerical-whole-section-fs,def:native-line-whole-section-general-fs,lem:native-numerical-line-tensor-power}
Every actual original fixed-line FS function for $L$ and rational
coefficient $a$ is numerical FS of its actual class $a[L]$, with the same
entire original function. Its given whole datum becomes an original
rational-line datum using its actual tensor-power line, its same positive
denominator, whole sections and original real weights.
\end{theorem}
\begin{proof}
If the original fixed-line datum has degree $d$ and denominator $m$, its
actual sections already belong to the whole line $L^{\otimes d}$.
Keep them, all weights and $m$ unchanged. Its true norm evaluation is
literally unchanged. The actual tensor-power class is $d[L]$. The original
degree relation is $d=ma$, hence its numerical class is
$m^{-1}d[L]=a[L]$. The original nonzero family remains the same family,
so it defines the claimed numerical FS function.
\end{proof}

\begin{remark}
These definitions use the actual original curve-degree numerical quotient
and true whole original sections, independently of any function-field
model comparison. They follow the varying rational-line shape of BJ's
\texttt{defi:FS} and \texttt{defi:psh}, used in FHX's psh preliminaries.
Rational-weight or ample-perturbation approximation requires its own
geometric theorem. In particular, this definition does not assert that a
fixed-line FS predicate descends through arbitrary numerical equivalence.
\end{remark}
